\documentclass[10pt,a4paper]{amsart}
\usepackage[margin=19mm]{geometry}
\usepackage{amsmath,amssymb,mathtools}
\usepackage[scr=rsfs,cal=euler]{mathalfa}
\usepackage{xfrac}
\usepackage[unicode,hidelinks,bookmarks=false]{hyperref}
\newcommand{\ie}{i.e.\ }
\newcommand{\cf}{cf.\ }
\newcommand{\resp}{respectively\ }
\newcommand{\Subsection}{\subsection}
\providecommand{\nobreakdash}{\nobreak\hbox{-}}
\setlength{\emergencystretch}{2em}
\multlinegap0pt
\usepackage{algorithm}\usepackage{algpseudocode}
\usepackage{amssymb}
\usepackage{float}
\newtheorem{lemma}{Lemma}
\newcommand{\ZZ}{\mathbb{Z}}
\title{\normalfont Sharp Structure-Agnostic Minimax Risk for Partial Linear Models}
\author{Haichen Hu, David Simchi-Levi}
\date{Source: arXiv:2609.07997v1 (CC BY 4.0)}
\begin{document}
\maketitle

\noindent\emph{Source and licence.} \normalfont Sharp Structure-Agnostic Minimax Risk for Partial Linear Models. Version arXiv:2609.07997v1 (CC BY 4.0), distributed by arXiv under the Creative Commons Attribution 4.0 International licence, http://creativecommons.org/licenses/by/4.0/, as recorded in the arXiv licence metadata for this exact version and shown on the abstract page https://arxiv.org/abs/2609.07997v1. Full licence notice: \texttt{licenses/arxiv-2609.07997-CC-BY-4.0.txt}.\par\smallskip

\noindent\textit{Editorial scope.} Counted unit: the article's lemma lem:chi2\_dist (source0.tex lines 1287--1296). The article's complete original proof of it is delivered with it (source0.tex lines 2606--2909, one \texttt{proof} environment of 304 lines, which the article places away from the statement). No mathematics is written, repaired or re-derived: the statement and the proof are the article's own lines, in the article's own notation, and no retained line is substituted or re-flowed.  No further source-local context is needed: every cross-reference of every retained line resolves inside the two delivered spans.  Nothing was omitted from any retained span, so the package carries no omission record.  No retained line contains a citation, so the wrapper carries no bibliography and no bibliography entry.  No retained line contains a figure environment, an image-inclusion command or drawing code, so no image is delivered and the package declares no asset dependency.  Editorial formatting only, in addition: an AMS \texttt{amsart} preamble that restates the article's own declarations and macros used by the retained lines, namely the environments \texttt{lemma} on separate counters, together with the standard packages those lines need.  The retained environments print in this document as Lemma 1 in order of appearance, whereas the article prints the same items with the numbering of its own full document.  The complete native source of the article as distributed in the arXiv e-print for this exact version is bundled unchanged as \texttt{sources/209683/source0.tex}.
\begin{lemma}\label{lem:chi2_dist}
    Set $D=\frac{\alpha^2+r^2-2b^2}{1-b^2}$. For any $N\in\ZZ^+$, we denote $\overline{P}_N$ as the uniform mixture of the \(M\) \(N\)-sample distributions \(P_j^{\otimes N}\), i.e., $\overline{P}_N=\frac{1}{M}\sum_{j=1}^MP_j^{\otimes N}$. Then we have that
    \[
    \chi^2(\overline{P}_N\|P_*^{\otimes N})=\frac{(1+D)^N-1}{M}.
    \]
    Hence, we have
    \[
    \texttt{D}_\text{TV}(\overline{P}_N\|P_*^{\otimes N})\le \frac{1}{2}\sqrt{\frac{(1+D)^N-1}{M}}.
    \]
\end{lemma}

\begin{proof}[Proof of Lemma \ref{lem:chi2_dist}]
Recall from the construction that \(b=\alpha r\).  Define the
one-observation likelihood ratios with respect to \(Q\) by
\[
L_*(x,t,y)
:=
\frac{dP_*}{dQ}(x,t,y)
=
1+bty
\]
and
\[
L_j(x,t,y)
:=
\frac{dP_j}{dQ}(x,t,y)
=
\{1+\alpha t\phi_j(x)\}
\{1+ry\phi_j(x)\},
\ 1\leq j\leq M.
\]
Because \(0\leq\alpha,r\leq1/8\), we have
\(0\leq b\leq1/64\), and hence,
\[
L_*(x,t,y)=1+bty\geq1-b>0.
\]
Thus, every \(P_j\) is absolutely continuous with respect to \(P_*\).

We first calculate the one-observation likelihood cross moments.  For
\(1\leq j,k\leq M\), define
\[
A_{jk}
:=
\mathbb E_{P_*}
\left[
\frac{dP_j}{dP_*}(X,T,Y)
\frac{dP_k}{dP_*}(X,T,Y)
\right].
\]
Since \(dP_*=L_*\,dQ\) and \(dP_j/dP_*=L_j/L_*\),
\begin{equation}
A_{jk}
=
\mathbb E_Q
\left[
\frac{L_j(X,T,Y)L_k(X,T,Y)}
     {L_*(X,T,Y)}
\right].
\label{eq:one-observation-cross-moment}
\end{equation}

Define $S:=\alpha T+rY,
\ 
U:=TY$. Because \(\phi_j(X)^2=1\) and \(b=\alpha r\), we have
\begin{align}
L_j(X,T,Y)
&=
\{1+\alpha T\phi_j(X)\}
\{1+rY\phi_j(X)\}
\nonumber\\
&=
1+\alpha T\phi_j(X)+rY\phi_j(X)
+\alpha rTY\phi_j(X)^2
\nonumber\\
&=
1+bTY+\phi_j(X)(\alpha T+rY)
\nonumber\\
&=
L_*(X,T,Y)+\phi_j(X)S.
\label{eq:likelihood-decomposition}
\end{align}
The same identity holds with \(j\) replaced by \(k\).  Therefore, we obtain
\begin{align}
\frac{L_jL_k}{L_*}
=
\frac{
\{L_*+\phi_j(X)S\}
\{L_*+\phi_k(X)S\}
}{L_*}
=
L_*
+\{\phi_j(X)+\phi_k(X)\}S
+\phi_j(X)\phi_k(X)\frac{S^2}{L_*}.
\label{eq:cross-integrand-expansion}
\end{align}

Under \(Q\), \(X\) is independent of \((T,Y)\).  Moreover, \(T\) and \(Y\)
are independent Rademacher random variables.  Consequently,
\[
\mathbb E_Q[L_*]
=
\mathbb E_Q[1+bTY]
=1,
\ 
\mathbb E_Q[S]
=
\alpha\mathbb E_Q[T]+r\mathbb E_Q[Y]
=0.
\]
Taking expectations in \eqref{eq:cross-integrand-expansion} and using the
independence of \(X\) and \((T,Y)\), we have
\begin{equation}
A_{jk}
=
1+
\mathbb E_X[\phi_j(X)\phi_k(X)]
\mathbb E_{T,Y}
\left[
\frac{S^2}{1+bU}
\right].
\label{eq:cross-moment-reduced}
\end{equation}

It remains to evaluate the second expectation.  Because \(T\) and \(Y\)
are independent Rademacher variables, \(U=TY\) is also Rademacher, so $\mathbb E[U]=0,
\ U^2=1$.
Moreover, notice that
\[
S^2
=
(\alpha T+rY)^2
=
\alpha^2+r^2+2\alpha rTY
=
\alpha^2+r^2+2bU,\ \frac1{1+bU}
=
\frac{1-bU}{1-b^2}.
\]
It follows that
\begin{align}
\mathbb E_{T,Y}
\left[
\frac{S^2}{1+bU}
\right]
&=
\frac1{1-b^2}
\mathbb E
\left[
\{\alpha^2+r^2+2bU\}(1-bU)
\right]
\nonumber\\
&=
\frac1{1-b^2}
\mathbb E
\left[
\alpha^2+r^2
+\{2b-b(\alpha^2+r^2)\}U
-2b^2U^2
\right]
\nonumber\\
&=
\frac{\alpha^2+r^2-2b^2}{1-b^2}
\nonumber\\
&=
D.
\label{eq:diagonal-increment}
\end{align}
Combining \eqref{eq:cross-moment-reduced} and
\eqref{eq:diagonal-increment}, we obtain that
\[
A_{jk}
=
1+D\mathbb E_X[\phi_j(X)\phi_k(X)].
\]
The code functions are orthonormal in \(L_2(\nu_d)\), so we have $\mathbb E_X[\phi_j(X)\phi_k(X)]
=
\mathbf 1\{j=k\}$.

Consequently,
\begin{equation}
A_{jk}
=
\begin{cases}
1+D, & j=k,\\
1,   & j\neq k.
\end{cases}
\label{eq:cross-moment-final}
\end{equation}
In particular, every off-diagonal cross moment is exactly one.

We now pass to \(N\) observations.  Write
\(Z_i=(X_i,T_i,Y_i)\).  By the product structure,
\[
\frac{dP_j^{\otimes N}}{dP_*^{\otimes N}}(Z_1,\ldots,Z_N)
=
\prod_{i=1}^N\frac{L_j(Z_i)}{L_*(Z_i)}.
\]
Hence,
\begin{equation}
\frac{d\overline P_N}{dP_*^{\otimes N}}(Z_1,\ldots,Z_N)
=
\frac1M\sum_{j=1}^M
\prod_{i=1}^N\frac{L_j(Z_i)}{L_*(Z_i)}.
\label{eq:mixture-likelihood-ratio}
\end{equation}
Recall that
\[
\chi^2(P\|Q)
:=
\mathbb E_Q
\left[
\left(\frac{dP}{dQ}-1\right)^2
\right]
=
\mathbb E_Q
\left[
\left(\frac{dP}{dQ}\right)^2
\right]-1.
\]
Using \eqref{eq:mixture-likelihood-ratio}, expanding the square, and using
independence under \(P_*^{\otimes N}\), we obtain
\begin{align}
1+
\chi^2\!\left(\overline P_N\middle\|P_*^{\otimes N}\right)
&=
\mathbb E_{P_*^{\otimes N}}
\left[
\left(
\frac{d\overline P_N}{dP_*^{\otimes N}}
\right)^2
\right] \nonumber\\
&=
\frac1{M^2}
\sum_{j=1}^M\sum_{k=1}^M
\mathbb E_{P_*^{\otimes N}}
\left[
\prod_{i=1}^N
\frac{L_j(Z_i)L_k(Z_i)}{L_*(Z_i)^2}
\right]
\nonumber\\
&=
\frac1{M^2}
\sum_{j=1}^M\sum_{k=1}^M
\left\{
\mathbb E_{P_*}
\left[
\frac{L_j(Z)L_k(Z)}{L_*(Z)^2}
\right]
\right\}^{N}
\nonumber\\
&=
\frac1{M^2}
\sum_{j=1}^M\sum_{k=1}^M A_{jk}^N.
\label{eq:mixture-second-moment}
\end{align}
By \eqref{eq:cross-moment-final}, the last double sum contains \(M\)
diagonal terms equal to \((1+D)^N\) and \(M(M-1)\) off-diagonal terms
equal to one.  Thus, we have
\begin{align*}
1+
\chi^2\!\left(\overline P_N\middle\|P_*^{\otimes N}\right)
&=
\frac{M(1+D)^N+M(M-1)}{M^2}
=
1+\frac{(1+D)^N-1}{M}.
\end{align*}
Subtracting one from both sides, we prove that
\[
\chi^2\!\left(\overline P_N\middle\|P_*^{\otimes N}\right)
=
\frac{(1+D)^N-1}{M}.
\]
Finally, by definition, we know that
\[
d_{\mathrm{TV}}(P,Q)
:=
\frac12
\mathbb E_Q
\left|
\frac{dP}{dQ}-1
\right|.
\]
By the Cauchy-Schwarz inequality, we have
\begin{align*}
d_{\mathrm{TV}}\!\left(\overline P_N,P_*^{\otimes N}\right)
&=
\frac12
\mathbb E_{P_*^{\otimes N}}
\left|
\frac{d\overline P_N}{dP_*^{\otimes N}}-1
\right|
\leq
\frac12
\left\{
\mathbb E_{P_*^{\otimes N}}
\left[
\left(
\frac{d\overline P_N}{dP_*^{\otimes N}}-1
\right)^2
\right]
\right\}^{1/2}
=
\frac12
\sqrt{
\chi^2\!\left(\overline P_N\middle\|P_*^{\otimes N}\right)
}
\end{align*}
Thus, we prove that
\[
d_{\mathrm{TV}}\!\left(\overline P_N,P_*^{\otimes N}\right)=
\frac12
\sqrt{\frac{(1+D)^N-1}{M}}.
\]
This proves the second claim.
\end{proof}

\end{document}
